\begin{textsample}{POS Dim 6 – ce – Score 9.00 – ce/ce\_ef\_25.txt}  \label{ex:f6_pos_037}
4.2.3.1.1.
O letramento científico no componente curricular Ciências O desenvolvimento científico e tecnológico da sociedade contemporânea imprime, ao currículo formal de Ciências, a necessidade de trabalhar em uma perspectiva para além da reprodução conceitual e experimental, tradicionalmente executadas nas escolas brasileiras.
O letramento científico se vincula a ideia de que é necessário agir sob o fundamento científico, ensejando processos e procedimentos próprios da ciência, isto é, empregar o conhecimento científico em diversas questões, considerando o processo investigativo.
Nesse aspecto, a Base Nacional Comum Curricular apresenta o conceito de letramento científico, enquanto compromisso do ensino de Ciências ao longo do Ensino Fundamental.
Para tanto, o documento explicita a necessidade da proposição, no currículo local, de oportunidades para o desenvolvimento da capacidade de compreender e interpretar o mundo, transformando -o a partir dos aportes teóricos advindos do contato com os objetos curriculares pertinentes a cada unidade temática.
Trata -se de proporcionar acesso ao repertório de conhecimentos, gerados historicamente pela humanidade, bem como fortalecer o processo do letramento linguístico, voltado para a educação científica, a qual assume um papel estratégico na construção de uma sociedade mais preparada para debruçar -se sobre seus próprios \textbf{problemas} e para formular as soluções viáveis.
Assim, o letramento científico compreende um conjunto de habilidades e saberes distribuídos em quatro dimensões : ( i ) o conhecimento do vocabulário, processos e produtos relativos à ciência e \textbf{tecnologia} ; ( ii ) a abertura mental para a descoberta, permitindo a (re)construção de saberes a partir da investigação e modelagem ; ( iii ) a comunicação com o uso correto do vocabulário e sistemas semióticos específicos das Ciências e, ( iv ) a intervenção na realidade pautadas nos princípios da ética e \textbf{sustentabilidade} socioambiental ( BRASIL, 2017 ).
Para que a cognição relacionada às quatro dimensões do letramento científico se \textbf{processe}, torna -se necessário que as/os professores de Ciências, ao desenvolverem o currículo escolar, o plano de ensino e, mais especificamente os planos de aula, promovam situações nas quais as/os estudantes possam interagir com os objetos específicos propostos na \textbf{matriz} do componente de Ciências, desenvolvendo as habilidades apresentadas no Quadro 02.
Quadro 02 : Dimensões do letramento científico. | Dimensão | Habilidades relacionadas ao letramento científico | |---|---| | Analisar demandas, delinear \textbf{problemas} e planejar investigações. | • Definir \textbf{problemas} do cotidiano ; • Observar o mundo a sua volta e fazer perguntas. • Propor hipóteses. | | Levantar hipóteses, analisar e representar semioticamente a realidade e os achados científicos. | • Planejar e realizar atividades de campo ( experimentos, observações, leituras, visitas, ambientes virtuais etc. ). • Avaliar informação ( validade, coerência e adequação ao \textbf{problema} formulado ). • Elaborar explicações e/ou modelos. • Associar explicações e/ou modelos à evolução histórica dos conhecimentos científicos envolvidos. • Aprimorar seus saberes e incorporar, gradualmente, e de modo significativo, o conhecimento científico. • Relatar informações de forma oral, escrita ou multimodal. • Participar de discussões de caráter científico com colegas, professores e professoras, familiares e comunidade em geral. • Desenvolver e utilizar ferramentas, inclusive \textbf{digitais}, para coleta, análise e representação de dados ( imagens, esquemas, tabelas, gráficos, quadros, diagramas, mapas, modelos, representações de sistemas, fluxogramas, mapas conceituais, simulações, \textbf{aplicativos} etc. ). | | Levantar hipóteses, analisar e representar semioticamente a realidade e os achados científicos. | • Selecionar e construir \textbf{argumentos} com base em evidências, modelos e/ou conhecimentos científicos. • Desenvolver soluções para \textbf{problemas} cotidianos usando diferentes ferramentas, inclusive \textbf{digitais}. • Apresentar, de forma sistemática, dados e resultados de investigações. • Considerar contra-argumentos para rever processos investigativos e conclusões. | | Intervir | • Desenvolver ações de intervenção para melhorar a qualidade de vida individual, coletiva e socioambiental. • Implementar soluções e avaliar sua eficácia para resolver \textbf{problemas} cotidianos. | Fonte : adaptado de BNCC, 2017.
Denotamos que o letramento científico está vinculado à inserção consciente na sociedade contemporânea, proporcionando às/aos jovens que adentrarão o ensino médio, o engajamento nos debates relativos ao uso consciente das \textbf{tecnologias} derivadas do desenvolvimento científico.
Para tanto, a educação científica para a cidadania deve movimentar saberes de todas as áreas do conhecimento humano, ampliando a capacidade das/dos estudantes de estabelecer as relações de causalidade em processos e \textbf{problemas} \textbf{ambientais} do meio no qual está inserido.
Esperamos, portanto, que, ao longo de sua trajetória escolar, as/os estudantes possam desenvolver o letramento científico, sendo capazes de explicar fatos cotidianos e propor ou elaborar soluções embasadas em conhecimentos científicos que foram (re)construídos a partir de sua reflexão crítica, o que envolve interligação entre aspectos cognitivos, experienciais e contextuais.
Conhecendo e compreendendo a realidade que o rodeia, acreditamos que será mais fácil para elas/eles intervirem, de maneira consciente, antecipando e prevendo quais são as consequências das atitudes humanas sobre o meio ambiente.

% matched lemmas: ambiental, aplicativo, argumento, digital, matriz, problema, processar, sustentabilidade, tecnologia
\end{textsample}
